% Options for packages loaded elsewhere
% Options for packages loaded elsewhere
\PassOptionsToPackage{unicode}{hyperref}
\PassOptionsToPackage{hyphens}{url}
\PassOptionsToPackage{dvipsnames,svgnames,x11names}{xcolor}
%
\documentclass[
  letterpaper,
  DIV=11,
  numbers=noendperiod]{scrartcl}
\usepackage{xcolor}
\usepackage{amsmath,amssymb}
\setcounter{secnumdepth}{-\maxdimen} % remove section numbering
\usepackage{iftex}
\ifPDFTeX
  \usepackage[T1]{fontenc}
  \usepackage[utf8]{inputenc}
  \usepackage{textcomp} % provide euro and other symbols
\else % if luatex or xetex
  \usepackage{unicode-math} % this also loads fontspec
  \defaultfontfeatures{Scale=MatchLowercase}
  \defaultfontfeatures[\rmfamily]{Ligatures=TeX,Scale=1}
\fi
\usepackage{lmodern}
\ifPDFTeX\else
  % xetex/luatex font selection
\fi
% Use upquote if available, for straight quotes in verbatim environments
\IfFileExists{upquote.sty}{\usepackage{upquote}}{}
\IfFileExists{microtype.sty}{% use microtype if available
  \usepackage[]{microtype}
  \UseMicrotypeSet[protrusion]{basicmath} % disable protrusion for tt fonts
}{}
\makeatletter
\@ifundefined{KOMAClassName}{% if non-KOMA class
  \IfFileExists{parskip.sty}{%
    \usepackage{parskip}
  }{% else
    \setlength{\parindent}{0pt}
    \setlength{\parskip}{6pt plus 2pt minus 1pt}}
}{% if KOMA class
  \KOMAoptions{parskip=half}}
\makeatother
% Make \paragraph and \subparagraph free-standing
\makeatletter
\ifx\paragraph\undefined\else
  \let\oldparagraph\paragraph
  \renewcommand{\paragraph}{
    \@ifstar
      \xxxParagraphStar
      \xxxParagraphNoStar
  }
  \newcommand{\xxxParagraphStar}[1]{\oldparagraph*{#1}\mbox{}}
  \newcommand{\xxxParagraphNoStar}[1]{\oldparagraph{#1}\mbox{}}
\fi
\ifx\subparagraph\undefined\else
  \let\oldsubparagraph\subparagraph
  \renewcommand{\subparagraph}{
    \@ifstar
      \xxxSubParagraphStar
      \xxxSubParagraphNoStar
  }
  \newcommand{\xxxSubParagraphStar}[1]{\oldsubparagraph*{#1}\mbox{}}
  \newcommand{\xxxSubParagraphNoStar}[1]{\oldsubparagraph{#1}\mbox{}}
\fi
\makeatother


\usepackage{longtable,booktabs,array}
\usepackage{calc} % for calculating minipage widths
% Correct order of tables after \paragraph or \subparagraph
\usepackage{etoolbox}
\makeatletter
\patchcmd\longtable{\par}{\if@noskipsec\mbox{}\fi\par}{}{}
\makeatother
% Allow footnotes in longtable head/foot
\IfFileExists{footnotehyper.sty}{\usepackage{footnotehyper}}{\usepackage{footnote}}
\makesavenoteenv{longtable}
\usepackage{graphicx}
\makeatletter
\newsavebox\pandoc@box
\newcommand*\pandocbounded[1]{% scales image to fit in text height/width
  \sbox\pandoc@box{#1}%
  \Gscale@div\@tempa{\textheight}{\dimexpr\ht\pandoc@box+\dp\pandoc@box\relax}%
  \Gscale@div\@tempb{\linewidth}{\wd\pandoc@box}%
  \ifdim\@tempb\p@<\@tempa\p@\let\@tempa\@tempb\fi% select the smaller of both
  \ifdim\@tempa\p@<\p@\scalebox{\@tempa}{\usebox\pandoc@box}%
  \else\usebox{\pandoc@box}%
  \fi%
}
% Set default figure placement to htbp
\def\fps@figure{htbp}
\makeatother





\setlength{\emergencystretch}{3em} % prevent overfull lines

\providecommand{\tightlist}{%
  \setlength{\itemsep}{0pt}\setlength{\parskip}{0pt}}



 


\KOMAoption{captions}{tableheading}
\makeatletter
\@ifpackageloaded{tcolorbox}{}{\usepackage[skins,breakable]{tcolorbox}}
\@ifpackageloaded{fontawesome5}{}{\usepackage{fontawesome5}}
\definecolor{quarto-callout-color}{HTML}{909090}
\definecolor{quarto-callout-note-color}{HTML}{0758E5}
\definecolor{quarto-callout-important-color}{HTML}{CC1914}
\definecolor{quarto-callout-warning-color}{HTML}{EB9113}
\definecolor{quarto-callout-tip-color}{HTML}{00A047}
\definecolor{quarto-callout-caution-color}{HTML}{FC5300}
\definecolor{quarto-callout-color-frame}{HTML}{acacac}
\definecolor{quarto-callout-note-color-frame}{HTML}{4582ec}
\definecolor{quarto-callout-important-color-frame}{HTML}{d9534f}
\definecolor{quarto-callout-warning-color-frame}{HTML}{f0ad4e}
\definecolor{quarto-callout-tip-color-frame}{HTML}{02b875}
\definecolor{quarto-callout-caution-color-frame}{HTML}{fd7e14}
\makeatother
\makeatletter
\@ifpackageloaded{caption}{}{\usepackage{caption}}
\AtBeginDocument{%
\ifdefined\contentsname
  \renewcommand*\contentsname{Table of contents}
\else
  \newcommand\contentsname{Table of contents}
\fi
\ifdefined\listfigurename
  \renewcommand*\listfigurename{List of Figures}
\else
  \newcommand\listfigurename{List of Figures}
\fi
\ifdefined\listtablename
  \renewcommand*\listtablename{List of Tables}
\else
  \newcommand\listtablename{List of Tables}
\fi
\ifdefined\figurename
  \renewcommand*\figurename{Figure}
\else
  \newcommand\figurename{Figure}
\fi
\ifdefined\tablename
  \renewcommand*\tablename{Table}
\else
  \newcommand\tablename{Table}
\fi
}
\@ifpackageloaded{float}{}{\usepackage{float}}
\floatstyle{ruled}
\@ifundefined{c@chapter}{\newfloat{codelisting}{h}{lop}}{\newfloat{codelisting}{h}{lop}[chapter]}
\floatname{codelisting}{Listing}
\newcommand*\listoflistings{\listof{codelisting}{List of Listings}}
\makeatother
\makeatletter
\makeatother
\makeatletter
\@ifpackageloaded{caption}{}{\usepackage{caption}}
\@ifpackageloaded{subcaption}{}{\usepackage{subcaption}}
\makeatother
\usepackage{bookmark}
\IfFileExists{xurl.sty}{\usepackage{xurl}}{} % add URL line breaks if available
\urlstyle{same}
\hypersetup{
  pdftitle={Lecture 1: Introduction},
  colorlinks=true,
  linkcolor={blue},
  filecolor={Maroon},
  citecolor={Blue},
  urlcolor={Blue},
  pdfcreator={LaTeX via pandoc}}


\title{Lecture 1: Introduction}
\author{}
\date{}
\begin{document}
\maketitle


\setcounter{secnumdepth}{3}

\section*{Learning Objectives}\label{learning-objectives}
\addcontentsline{toc}{section}{Learning Objectives}

After reading this lecture you should be able to:

\begin{itemize}
\tightlist
\item
  \textbf{Describe the fundamental features of bonds.}
\item
  \textbf{Explain the major types of bond issuers and the sectors of the
  bond market.}
\item
  \textbf{Interpret the importance of the term to maturity in bond
  valuation and risk.}
\item
  \textbf{Differentiate between fixed‑rate, floating‑rate and
  inverse‑floating‑rate securities.}
\item
  \textbf{Explain what embedded options are and how options affect a
  bond's cash flows.}
\item
  \textbf{Identify and describe the common types of embedded options
  (call, put, conversion, currency choice).}
\item
  \textbf{Discuss convertible bonds and exchangeable bonds.}
\item
  \textbf{Recognize the spectrum of risks faced by fixed‑income
  investors.}
\item
  \textbf{Describe the role and functioning of the secondary bond
  market.}
\end{itemize}

These objectives come directly from the first chapter of Fabozzi's text
and frame the topics covered in this lecture.

\begin{center}\rule{0.5\linewidth}{0.5pt}\end{center}

\section{What is a bond?}\label{what-is-a-bond}

A \textbf{bond} is a debt instrument requiring the issuer (the borrower)
to repay the amount borrowed plus interest over a specified period of
time. In the simplest ``plain vanilla'' bond the issuer promises a fixed
series of cash flows: periodic coupon payments and the repayment of
principal at maturity. Because the cash flow pattern is known up front
(assuming the issuer does not default or redeem the bond early), an
investor who holds the bond until the maturity date is assured of
receiving a predetermined sequence of payments.

\subsection{Bond cash‑flow structure}\label{bond-cashflow-structure}

Two contractual elements determine a bond's cash flow pattern:

\begin{enumerate}
\def\labelenumi{\arabic{enumi}.}
\tightlist
\item
  \textbf{Principal (face value or par value)} -- the amount the issuer
  must repay at the maturity date. Synonyms include redemption value,
  maturity value and par value.\\
\item
  \textbf{Coupon rate and payments} -- the interest rate the issuer
  agrees to pay each year. The coupon rate multiplied by the principal
  gives the annual coupon payment. For example, a bond with a 6\,\%
  coupon rate and a face value of \$1\,000 pays \$60 in interest each
  year, usually in two semi‑annual installments.
\end{enumerate}

The maturity date is the date when the principal is repaid. Traditional
fixed‑rate bonds promise equal coupon payments over the life of the
bond. Zero‑coupon bonds make no periodic interest payments; instead they
are sold at a discount and repay the principal at maturity.

\begin{tcolorbox}[enhanced jigsaw, titlerule=0mm, rightrule=.15mm, bottomtitle=1mm, colframe=quarto-callout-note-color-frame, opacitybacktitle=0.6, coltitle=black, breakable, toprule=.15mm, title=\textcolor{quarto-callout-note-color}{\faInfo}\hspace{0.5em}{Worked Example~1~--~Coupon calculation}, leftrule=.75mm, colbacktitle=quarto-callout-note-color!10!white, left=2mm, arc=.35mm, opacityback=0, toptitle=1mm, bottomrule=.15mm, colback=white]

Suppose a bond has a principal of \$1\,000 and a coupon rate of 5\,\%.
The annual coupon payment is simply (0.05×1\,000=50). If coupons are
paid semi‑annually, the investor receives \$25 every six months. This
illustrates how the coupon rate and principal determine the bond's cash
flow pattern.

\end{tcolorbox}

\subsection{Term to maturity}\label{term-to-maturity}

The \textbf{term to maturity} is the number of years over which the
issuer must meet its obligations. Bonds with maturity of 1--5~years are
called \textbf{short‑term}, those with 5--12~years are
\textbf{intermediate‑term}, and bonds with more than 12~years are
\textbf{long‑term}. Maturity matters because:

\begin{itemize}
\tightlist
\item
  It determines the horizon over which the investor receives coupon
  payments and principal.
\item
  The yield on a bond depends on its maturity; yields for different
  maturities form the \textbf{yield curve}.
\item
  Price volatility increases with maturity: all else equal, long‑term
  bond prices fluctuate more when yields change.
\end{itemize}

\subsection{Principal and coupon
types}\label{principal-and-coupon-types}

Most bonds have a \textbf{fixed coupon rate}; however the bond market
also includes \textbf{floating‑rate bonds} in which the coupon resets
periodically based on a reference rate plus a quoted margin. For
example, a floating‑rate bond might pay 1‑month LIBOR plus 150\,basis
points (1.5\,\%). LIBOR is the London Interbank Offered Rate, an
interbank lending rate reported for several currencies. If 1‑month LIBOR
is 3.5\,\% and the quoted margin is 1.5\,\%, the coupon for that period
resets at 5.0\,\%.

\begin{tcolorbox}[enhanced jigsaw, titlerule=0mm, rightrule=.15mm, bottomtitle=1mm, colframe=quarto-callout-note-color-frame, opacitybacktitle=0.6, coltitle=black, breakable, toprule=.15mm, title=\textcolor{quarto-callout-note-color}{\faInfo}\hspace{0.5em}{Market update -- LIBOR transition}, leftrule=.75mm, colbacktitle=quarto-callout-note-color!10!white, left=2mm, arc=.35mm, opacityback=0, toptitle=1mm, bottomrule=.15mm, colback=white]

LIBOR is no longer the standard benchmark for newly issued floating‑rate
instruments. After concerns about manipulation and the decline of
unsecured interbank lending activity, regulators and market participants
moved to alternative reference rates. U.S. dollar LIBOR has largely been
replaced by \textbf{SOFR} (Secured Overnight Financing Rate), which is
based on transactions in the U.S. Treasury repo market. In Canada,
Canadian dollar LIBOR was not the main local benchmark; the key
transition has been from \textbf{CDOR} to \textbf{CORRA} (Canadian
Overnight Repo Rate Average). In other major markets, common
replacements include \textbf{SONIA} in the United Kingdom, \textbf{€STR}
in the euro area, \textbf{SARON} in Switzerland and \textbf{TONA} in
Japan.

\end{tcolorbox}

\begin{tcolorbox}[enhanced jigsaw, titlerule=0mm, rightrule=.15mm, bottomtitle=1mm, colframe=quarto-callout-note-color-frame, opacitybacktitle=0.6, coltitle=black, breakable, toprule=.15mm, title=\textcolor{quarto-callout-note-color}{\faInfo}\hspace{0.5em}{What is the repo market?}, leftrule=.75mm, colbacktitle=quarto-callout-note-color!10!white, left=2mm, arc=.35mm, opacityback=0, toptitle=1mm, bottomrule=.15mm, colback=white]

A \textbf{repurchase agreement}, or \textbf{repo}, is a short‑term
secured loan. One party sells a security, often a U.S. Treasury, for
cash and agrees to buy it back later at a slightly higher price.
Economically, the security seller is borrowing cash and pledging the
security as collateral; the difference between the sale price and
repurchase price implies the \textbf{repo rate}. Because SOFR is based
on overnight Treasury repo transactions, it reflects the cost of
borrowing cash secured by Treasury collateral rather than unsecured bank
borrowing.

\includegraphics[width=7.41in,height=2.29in]{lec1_files/figure-latex/mermaid-figure-1.png}

In a \textbf{reverse repo}, the same transaction is viewed from the cash
lender's side: the lender provides cash, receives collateral and later
returns the collateral when repaid.

\end{tcolorbox}

\textbf{Inverse floating‑rate bonds} move in the opposite direction: the
coupon falls when reference rates rise and vice versa. Deferred‑coupon
bonds allow the issuer to avoid cash outflows for a specified period;
examples include deferred‑interest bonds, step‑up bonds and
payment‑in‑kind bonds. These structures were used to accommodate
leveraged buyouts in the 1980s.

\begin{tcolorbox}[enhanced jigsaw, titlerule=0mm, rightrule=.15mm, bottomtitle=1mm, colframe=quarto-callout-note-color-frame, opacitybacktitle=0.6, coltitle=black, breakable, toprule=.15mm, title=\textcolor{quarto-callout-note-color}{\faInfo}\hspace{0.5em}{Inflation-Linked Bonds: TIPS and Linkers}, leftrule=.75mm, colbacktitle=quarto-callout-note-color!10!white, left=2mm, arc=.35mm, opacityback=0, toptitle=1mm, bottomrule=.15mm, colback=white]

\textbf{Inflation-linked bonds}, also known as ``linkers,'' are bonds
whose principal and/or coupon payments adjust with changes in an
official inflation index. The most prominent examples are
\textbf{Treasury Inflation-Protected Securities (TIPS)} in the United
States, which are indexed to the Consumer Price Index (CPI).

\begin{itemize}
\tightlist
\item
  With TIPS, the principal amount is adjusted by the change in CPI;
  coupon payments are calculated on the inflation-adjusted principal, so
  both the interest payments and maturity value rise with inflation.
  This provides investors with protection against the eroding effect of
  inflation on fixed income.
\item
  In Canada, the equivalent instrument is called a \textbf{Real Return
  Bond (RRB)}, which adjusts both principal and coupon payments based on
  changes in the Canadian Consumer Price Index (CPI).
\item
  At maturity, investors in TIPS receive the greater of the original
  principal or the inflation-adjusted principal, guaranteeing repayment
  of their initial investment even in a period of deflation.
\end{itemize}

Inflation-linked bonds are an important segment of the global bond
market, providing a direct way to hedge inflation risk.

\end{tcolorbox}

\subsection{Amortization feature}\label{amortization-feature}

In a \textbf{non‑amortizing bond} the entire principal is repaid at
maturity. \textbf{Amortizing securities}, by contrast, repay principal
over time according to an amortization schedule. Mortgage‑backed and
asset‑backed securities are examples. For amortizing bonds investors
often compute the \textbf{weighted average life} to measure the
effective maturity.

Suppose an investor purchases a \textbf{5-year amortizing bond} with a
principal of \$1,000 and equal annual payments of principal and
interest. Unlike a non-amortizing bond---which pays all principal at
maturity---this bond returns part of the principal each year along with
interest.

\begin{longtable}[]{@{}lll@{}}
\toprule\noalign{}
Year & Principal Repaid & Remaining Principal \\
\midrule\noalign{}
\endhead
\bottomrule\noalign{}
\endlastfoot
1 & \$200 & \$800 \\
2 & \$200 & \$600 \\
3 & \$200 & \$400 \\
4 & \$200 & \$200 \\
5 & \$200 & \$0 \\
\end{longtable}

\textbf{Weighted Average Life (WAL):} Weighted average life measures the
average time until each principal dollar is repaid, weighted by the
amount of principal paid at each period. It is useful for comparing the
effective maturity of amortizing securities versus a standard bullet
(non-amortizing) bond.

To calculate WAL:

\[
\text{WAL} = \frac{\sum_{t=1}^{N} P_t \times t}{\text{Total Principal}}
\]

where \(P_t\) is the principal repaid at time \(t\).

For this bond:

\[
\text{WAL} = \frac{200 \times 1 + 200 \times 2 + 200 \times 3 + 200 \times 4 + 200 \times 5}{1000} = \frac{200 (1+2+3+4+5)}{1000} = \frac{200 \times 15}{1000} = \frac{3000}{1000} = 3\ \text{years}
\]

Thus, although the \textbf{final maturity} is 5 years, the
\textbf{weighted average life} is only 3 years, since principal is
returned gradually rather than all at once.

\subsection{Embedded options}\label{embedded-options}

Many bonds include options that allow either the issuer or the
bondholder to alter the cash flow pattern. Embedded options complicate
valuation because the timing and amount of cash flows depend on future
interest rates. Common options include:

\begin{itemize}
\tightlist
\item
  \textbf{Call provision} -- The issuer may redeem the bond before
  maturity at a specified price. Issuers call bonds when market rates
  fall, disadvantaging investors by limiting price appreciation and
  forcing reinvestment at lower rates.\\
\item
  \textbf{Put provision} -- The bondholder may sell the bond back to the
  issuer at par on specified dates. This protects investors when rates
  rise.
\item
  \textbf{Conversion option} -- Convertible bonds allow the bondholder
  to exchange the bond for a specified number of shares of common stock.
\item
  \textbf{Currency‑choice option} -- Payments may be made in different
  currencies, giving one party the ability to benefit from exchange‑rate
  movements.
\end{itemize}

Some securities contain multiple embedded options (e.g., a callable,
putable, convertible bond). These features must be valued using
option‑pricing techniques, which are covered in later lectures.

\begin{tcolorbox}[enhanced jigsaw, titlerule=0mm, rightrule=.15mm, bottomtitle=1mm, colframe=quarto-callout-note-color-frame, opacitybacktitle=0.6, coltitle=black, breakable, toprule=.15mm, title=\textcolor{quarto-callout-note-color}{\faInfo}\hspace{0.5em}{Convertible and exchangeable bonds}, leftrule=.75mm, colbacktitle=quarto-callout-note-color!10!white, left=2mm, arc=.35mm, opacityback=0, toptitle=1mm, bottomrule=.15mm, colback=white]

A \textbf{convertible bond} allows the bondholder to exchange the bond
for a predetermined number of the issuer's common shares. An
\textbf{exchangeable bond} allows exchange for shares of a company other
than the issuer. These bonds provide upside exposure to stock price
appreciation while retaining bond‑like features.

\end{tcolorbox}

\subsection{Identifying bond issues}\label{identifying-bond-issues}

With hundreds of thousands of bonds outstanding, securities need unique
identifiers. In the U.S. the \textbf{CUSIP} (Committee on Uniform
Security Identification Procedures) assigns a nine‑character code
identifying the issuer and issue. Debt issues covered include
Treasuries, municipals, corporates, mortgage‑backed, asset‑backed and
commercial paper. Market participants often identify a bond by its
issuer, coupon and maturity (e.g., Alcoa 5.95s 2/1/2037).

\section{Innovation and sectors of the bond
market}\label{innovation-and-sectors-of-the-bond-market}

\subsection{Financial innovation}\label{financial-innovation}

Beginning in the early 1980s a wide range of new bond structures was
introduced. The residential mortgage market saw the development of
\textbf{mortgage pass‑through securities}, where pools of individual
mortgages are securitized and cash flows are passed through to
investors. From these basic instruments issuers created derivative
securities such as \textbf{collateralized mortgage obligations (CMOs)}
and \textbf{stripped mortgage‑backed securities} to meet the investment
needs of a broader range of institutional investors.

\subsection{Sectors of the U.S. bond
market}\label{sectors-of-the-u.s.-bond-market}

The U.S. bond market -- the world's largest -- is divided into six
primary sectors:

\begin{longtable}[]{@{}
  >{\raggedright\arraybackslash}p{(\linewidth - 2\tabcolsep) * \real{0.3158}}
  >{\raggedright\arraybackslash}p{(\linewidth - 2\tabcolsep) * \real{0.6842}}@{}}
\toprule\noalign{}
\begin{minipage}[b]{\linewidth}\raggedright
Sector
\end{minipage} & \begin{minipage}[b]{\linewidth}\raggedright
Description
\end{minipage} \\
\midrule\noalign{}
\endhead
\bottomrule\noalign{}
\endlastfoot
\textbf{U.S. Treasury} & Bonds, notes and bills issued by the U.S.
government. Treasuries serve as benchmarks for valuation and interest
rates globally. \\
\textbf{Agency} & Debenture securities issued by federally related
institutions and government‑sponsored enterprises. These are unsecured
and constitute the smallest sector. \\
\textbf{Municipal} & Bonds issued by states, cities and local
authorities. Divided into tax‑exempt and taxable subsectors. Structures
include \textbf{tax‑backed bonds} and \textbf{revenue bonds}. \\
\textbf{Corporate} & Debt issued by U.S. corporations and non‑U.S.
corporations in the U.S. market. Instruments include bonds, medium‑term
notes, structured notes and commercial paper. Corporates are categorized
as investment‑grade or non‑investment‑grade based on credit ratings. \\
\textbf{Asset‑backed securities (ABS)} & Securities backed by pools of
loans or receivables such as auto loans, credit‑card receivables and
music royalties. Issuers include captive finance subsidiaries of
corporations; unusual ABS include ``Bowie bonds'' backed by music
royalties. \\
\textbf{Mortgage} & Securities backed by mortgage loans on residential
or commercial properties. The mortgage sector is subdivided into
residential and commercial mortgages, with further distinctions between
\textbf{agency} and \textbf{non‑agency} mortgage‑backed securities and
between loans to \textbf{prime} and \textbf{subprime} borrowers. \\
\end{longtable}

These sectors span a wide range of issuers and structures. For instance,
the municipal sector splits into tax‑exempt and taxable subsectors, with
tax‑exempt bonds historically dominating. The corporate sector is
divided into investment‑grade and high‑yield (non‑investment‑grade)
segments based on third‑party credit ratings. In the mortgage sector,
agency mortgage‑backed securities are backed by Ginnie Mae, Fannie Mae
or Freddie Mac, whereas non‑agency securities are backed by private
issuers.

\subsubsection{Asset‑backed and mortgage‑backed
markets}\label{assetbacked-and-mortgagebacked-markets}

The \textbf{asset‑backed securities sector} pools loans or receivables
and uses the asset pool as collateral. Captive finance companies
(subsidiaries of manufacturers that finance customer purchases) are
common issuers; examples include Ford Motor Credit and Caterpillar
Financial. Even performing artists have issued ABS backed by future
royalties.

The \textbf{residential mortgage sector} includes loans for one‑ to
four‑family homes. Loans are classified by borrower credit quality
(prime or subprime) and by whether they conform to underwriting
standards of a federal agency or government‑sponsored enterprise. Agency
mortgage‑backed securities (MBS) are backed by Ginnie Mae, Fannie Mae or
Freddie Mac; non‑agency MBS are issued by private entities and further
subdivided into prime and subprime sectors.

The \textbf{commercial mortgage sector} backs securities with loans on
income‑producing properties such as office buildings, shopping malls,
hotels and healthcare facilities.

\subsubsection{Pooled investment
vehicles}\label{pooled-investment-vehicles}

Retail and institutional investors can gain exposure to the bond market
through \textbf{pooled investment vehicles} such as investment company
shares, exchange‑traded funds, hedge funds and real estate investment
trusts. These vehicles offer diversification, liquidity and professional
management, making them an attractive alternative to purchasing
individual bonds.

\begin{center}\rule{0.5\linewidth}{0.5pt}\end{center}

\section{Risks faced by bond
investors}\label{risks-faced-by-bond-investors}

Fixed‑income investing involves numerous risks. Fabozzi lists nine key
risks. Understanding these risks is essential for effective portfolio
construction and risk control.

\subsection{Interest‑rate risk}\label{interestrate-risk}

Bond prices move inversely with changes in market interest rates: when
rates rise, bond prices fall, and vice versa. If an investor must sell
before maturity, rising rates may lead to capital losses. The price
sensitivity to rate changes depends on factors such as maturity, coupon
and embedded options.

\subsection{Reinvestment risk}\label{reinvestment-risk}

Yield calculations assume that coupon payments are reinvested at the
same rate. \textbf{Reinvestment risk} is the risk that the actual
reinvestment rate will be lower than expected. Reinvestment risk is
greater for high‑coupon bonds and longer holding periods. Interest‑rate
risk and reinvestment risk offset each other; strategies such as
immunization exploit this offset.

\subsection{Call risk}\label{call-risk}

A call provision allows the issuer to retire (call) the bond before
maturity. Investors face three disadvantages:

\begin{enumerate}
\def\labelenumi{\arabic{enumi}.}
\tightlist
\item
  Uncertainty in cash flow timing because the call date is unknown.\\
\item
  Increased reinvestment risk: bonds are usually called when rates
  decline, forcing investors to reinvest at lower yields.\\
\item
  Limited price appreciation: the bond's price may not rise above the
  call price, reducing capital gains potential.
\end{enumerate}

\subsection{Credit risk}\label{credit-risk}

\textbf{Credit risk} refers to the possibility that the issuer fails to
meet its contractual obligations. It has three components:

\begin{itemize}
\tightlist
\item
  \textbf{Default risk} -- the risk of missed interest or principal
  payments.\\
\item
  \textbf{Credit spread risk} -- the risk that the yield spread between
  a corporate bond and a comparable Treasury widens, causing the bond's
  price to decline.\\
\item
  \textbf{Downgrade risk} -- the risk that a rating agency lowers the
  bond's credit rating, which typically leads to a price decline.
\end{itemize}

Credit ratings assigned by agencies (e.g., Moody's, S\&P, Fitch) help
investors gauge default risk. Investors are also exposed to
\textbf{counterparty risk}, the risk that the counterparty in a
transaction fails to meet its obligations. For example, leveraged
purchases of bonds or trades in over‑the‑counter derivatives expose
investors to counterparty default risk.

\subsection{Inflation risk}\label{inflation-risk}

The real value of future cash flows depends on inflation.
\textbf{Inflation risk} arises when the purchasing power of a bond's
cash flows declines due to rising prices. Floating‑rate bonds tied to
inflation (also called \textbf{linkers}) reduce inflation risk.

\subsection{Exchange‑rate risk}\label{exchangerate-risk}

For bonds denominated in a foreign currency, the investor's return
depends on exchange rates. If the foreign currency depreciates relative
to the investor's home currency, the value of the cash flows declines.

\subsection{Liquidity risk}\label{liquidity-risk}

\textbf{Liquidity risk} reflects the ease with which a bond can be sold
at a reasonable price. A wider bid--ask spread signals higher liquidity
risk. While retail investors planning to hold to maturity may ignore
liquidity risk, institutional investors must mark portfolios to market
and therefore require market prices based on sufficient trading volume.

\subsection{Volatility risk}\label{volatility-risk}

For bonds with embedded options, price depends on both interest rates
and the expected volatility of rates. Greater expected volatility
increases the value of the option. For callable bonds and
mortgage‑backed securities the price tends to decrease when volatility
rises because the option given to the issuer becomes more valuable.
Volatility risk is often quantified using \textbf{vega} from option
pricing theory.

\subsection{Risk risk}\label{risk-risk}

Fabozzi introduces \textbf{risk risk}, defined as the risk of not
understanding the risks inherent in a security. Investors mitigate risk
risk by staying informed about analytical techniques and avoiding
securities they do not understand. Although complex securities may offer
higher returns, thorough analysis is essential.

\begin{center}\rule{0.5\linewidth}{0.5pt}\end{center}

\section{Secondary market for bonds}\label{secondary-market-for-bonds}

After a bond is issued in the \textbf{primary market}, it can be bought
and sold in the \textbf{secondary market}. Secondary trading provides
liquidity and allows investors to adjust portfolios in response to
changing market conditions. Liquidity varies across sectors: Treasury
securities trade heavily, while certain municipals or high‑yield
corporates may trade infrequently. The presence of active secondary
markets is crucial for price discovery and portfolio valuation.

\begin{center}\rule{0.5\linewidth}{0.5pt}\end{center}

\section*{Additional worked examples}\label{additional-worked-examples}
\addcontentsline{toc}{section}{Additional worked examples}

\begin{tcolorbox}[enhanced jigsaw, titlerule=0mm, rightrule=.15mm, bottomtitle=1mm, colframe=quarto-callout-note-color-frame, opacitybacktitle=0.6, coltitle=black, breakable, toprule=.15mm, title=\textcolor{quarto-callout-note-color}{\faInfo}\hspace{0.5em}{Worked Example~3~--~Pricing a zero‑coupon bond}, leftrule=.75mm, colbacktitle=quarto-callout-note-color!10!white, left=2mm, arc=.35mm, opacityback=0, toptitle=1mm, bottomrule=.15mm, colback=white]

Consider a zero‑coupon bond with a face value of \$1\,000 and 3~years to
maturity. The yield (discount rate) is 4~\%. Because there are no coupon
payments, the price is simply the present value of the principal:

\[
P = \frac{F}{(1+r)^T} = \frac{1000}{(1.04)^3} \approx \$889.00 .
\]

This example illustrates the use of the bond pricing formula for
zero‑coupon securities.

\end{tcolorbox}

\begin{tcolorbox}[enhanced jigsaw, titlerule=0mm, rightrule=.15mm, bottomtitle=1mm, colframe=quarto-callout-note-color-frame, opacitybacktitle=0.6, coltitle=black, breakable, toprule=.15mm, title=\textcolor{quarto-callout-note-color}{\faInfo}\hspace{0.5em}{Worked Example~4~--~Floating‑rate coupon reset}, leftrule=.75mm, colbacktitle=quarto-callout-note-color!10!white, left=2mm, arc=.35mm, opacityback=0, toptitle=1mm, bottomrule=.15mm, colback=white]

A floating‑rate bond pays \textbf{1‑month LIBOR + 150~basis points}. On
the reset date LIBOR is 3.75~\%. The coupon for the next period resets
to:

\[
\text{Coupon} = 3.75\% + 1.50\% = 5.25\% .
\]

If the bond's face value is \$1\,000 and coupons are paid quarterly, the
next coupon payment will be \[
\frac{5.25\%}{4} \times 1000 = 13.125 .
\]

This demonstrates how the reference rate and quoted margin determine
floating‑rate coupons.

\end{tcolorbox}

\begin{tcolorbox}[enhanced jigsaw, titlerule=0mm, rightrule=.15mm, bottomtitle=1mm, colframe=quarto-callout-note-color-frame, opacitybacktitle=0.6, coltitle=black, breakable, toprule=.15mm, title=\textcolor{quarto-callout-note-color}{\faInfo}\hspace{0.5em}{Worked Example~5~--~Evaluating call risk}, leftrule=.75mm, colbacktitle=quarto-callout-note-color!10!white, left=2mm, arc=.35mm, opacityback=0, toptitle=1mm, bottomrule=.15mm, colback=white]

Suppose a callable bond can be redeemed at \$1\,050 after two years. The
bond has a coupon rate of 6~\% and face value \$1\,000. If market
interest rates fall to 4~\%, the price of a comparable noncallable bond
may rise to \$1\,100. However, the issuer is likely to call the bond at
\$1\,050 to refinance at the lower rate. Investors therefore cannot
benefit fully from price appreciation and face reinvestment at lower
yields.

\end{tcolorbox}

\begin{tcolorbox}[enhanced jigsaw, titlerule=0mm, rightrule=.15mm, bottomtitle=1mm, colframe=quarto-callout-note-color-frame, opacitybacktitle=0.6, coltitle=black, breakable, toprule=.15mm, title=\textcolor{quarto-callout-note-color}{\faInfo}\hspace{0.5em}{Worked Example~6~--~Credit spread risk}, leftrule=.75mm, colbacktitle=quarto-callout-note-color!10!white, left=2mm, arc=.35mm, opacityback=0, toptitle=1mm, bottomrule=.15mm, colback=white]

Consider a corporate bond yielding 5~\% when a comparable Treasury
yields 3~\%. The credit spread is 2~\% (200~basis points). If economic
conditions deteriorate, the spread may widen to 3~\%, causing the
corporate bond's price to fall. Investors exposed to credit spread risk
must monitor factors affecting the issuer's credit quality.

\end{tcolorbox}

\begin{center}\rule{0.5\linewidth}{0.5pt}\end{center}

\section*{Key Points}\label{key-points}
\addcontentsline{toc}{section}{Key Points}

\begin{itemize}
\item
  A \textbf{bond} is a debt contract in which the issuer promises a
  schedule of cash flows---typically periodic \textbf{coupon} payments
  and repayment of \textbf{principal} at \textbf{maturity}---unless
  default or embedded options alter those flows.
\item
  \textbf{Principal (face or par value)} and the \textbf{coupon rate}
  determine the promised cash flows; traditional fixed‑rate bonds pay
  equal coupons, while \textbf{zero‑coupon} bonds pay interest
  implicitly through a discount to par.
\item
  \textbf{Term to maturity} shapes interest‑rate risk and yield
  quotation; yields by maturity trace the \textbf{yield curve}, and
  longer maturities generally imply greater \textbf{price volatility}
  when yields change.
\item
  Markets also include \textbf{floating‑rate} bonds (coupon resets off a
  reference rate plus a \textbf{quoted margin}), \textbf{inverse
  floating‑rate} bonds, deferred‑coupon structures, and
  \textbf{inflation‑linked} bonds such as U.S. \textbf{TIPS} and
  Canadian \textbf{RRBs}.
\item
  New floating‑rate issues increasingly reference \textbf{SOFR},
  \textbf{CORRA}, \textbf{SONIA}, \textbf{€STR}, and similar benchmarks
  rather than \textbf{LIBOR}.
\item
  \textbf{Non‑amortizing} bonds repay principal at maturity;
  \textbf{amortizing} securities repay principal over a schedule;
  \textbf{weighted average life (WAL)} summarizes how quickly principal
  is returned on amortizing structures.
\item
  \textbf{Embedded options}---\textbf{call}, \textbf{put},
  \textbf{conversion}, and \textbf{currency choice}---can change
  cash‑flow timing or upside and require option‑aware valuation.
\item
  \textbf{Convertible} bonds exchange into the issuer's equity;
  \textbf{exchangeable} bonds exchange into shares of another company.
\item
  Major risks include \textbf{interest‑rate risk}, \textbf{reinvestment
  risk}, \textbf{call risk}, \textbf{credit risk} (default,
  \textbf{credit spread}, downgrade), \textbf{counterparty risk},
  \textbf{inflation risk}, \textbf{exchange‑rate risk},
  \textbf{liquidity risk}, \textbf{volatility risk} for option‑rich
  bonds, and \textbf{risk risk}---investing in securities one does not
  fully understand.
\item
  Bonds are sold first in the \textbf{primary market} and then trade in
  the \textbf{secondary market}, where liquidity and price discovery
  vary widely by sector and issue.
\end{itemize}




\end{document}
